\documentclass[tikz, border=0.2mm]{standalone}
\usepackage{amsmath,xparse}
\newcommand{\tikzsetnextfilename}[1]{}


\ExplSyntaxOn
% 1. Create a variable to hold the path construction
\tl_new:N \l_lattice_path_tl
% 2. Define the main command
\NewDocumentCommand{\latticePath}{ O{} m m }
 {
  % Clear the variable for the new path
  \tl_clear:N \l_lattice_path_tl
  % Loop through the input string (#3)
  \str_map_inline:nn { #3 }
   {
    % Append the TikZ code to the variable based on the character
    \str_case:nnF { ##1 }
     {
       % Numpad
       {1}{ \tl_put_right:Nn \l_lattice_path_tl { --~++(-1,-1) } } % SW
       {2}{ \tl_put_right:Nn \l_lattice_path_tl { --~++(0,-1)  } } % S
       {3}{ \tl_put_right:Nn \l_lattice_path_tl { --~++(1,-1)  } } % SE
       {4}{ \tl_put_right:Nn \l_lattice_path_tl { --~++(-1,0)  } } % W
       {6}{ \tl_put_right:Nn \l_lattice_path_tl { --~++(1,0)   } } % E
       {7}{ \tl_put_right:Nn \l_lattice_path_tl { --~++(-1,1)  } } % NW
       {8}{ \tl_put_right:Nn \l_lattice_path_tl { --~++(0,1)   } } % N
       {9}{ \tl_put_right:Nn \l_lattice_path_tl { --~++(1,1)   } } % NE
       % Cardinals
       {n}{ \tl_put_right:Nn \l_lattice_path_tl { --~++(0,1)   } }
       {s}{ \tl_put_right:Nn \l_lattice_path_tl { --~++(0,-1)  } }
       {e}{ \tl_put_right:Nn \l_lattice_path_tl { --~++(1,0)   } }
       {w}{ \tl_put_right:Nn \l_lattice_path_tl { --~++(-1,0)  } }
       % Relatives
       {u}{ \tl_put_right:Nn \l_lattice_path_tl { --~++(0,1)   } }
       {d}{ \tl_put_right:Nn \l_lattice_path_tl { --~++(0,-1)  } }
       {l}{ \tl_put_right:Nn \l_lattice_path_tl { --~++(-1,0)  } }
       {r}{ \tl_put_right:Nn \l_lattice_path_tl { --~++(1,0)   } }
     }
     {
       % Fallback (ignore unknown characters)
     }
   }
   % 3. Execute the Draw Command using the constructed path
   % Note: We use ~ for spaces inside ExplSyntax
   \draw[line~width=4.0pt, line~cap=round, line~join=round, #1] (#2) \l_lattice_path_tl ;
 }
\ExplSyntaxOff


\begin{document}
\tikzsetnextfilename{dyckPath000}
\begin{tikzpicture}
    \draw[gray!30, thin] (0,0) grid (3,3);
    \latticePath{0,0}{nenene}
\end{tikzpicture}

\tikzsetnextfilename{dyckPath010}
\begin{tikzpicture}
    \draw[gray!30, thin] (0,0) grid (3,3);
    \latticePath{0,0}{nennee}
\end{tikzpicture}

\tikzsetnextfilename{dyckPath100}
\begin{tikzpicture}
    \draw[gray!30, thin] (0,0) grid (3,3);
    \latticePath{0,0}{nneene}
\end{tikzpicture}

\tikzsetnextfilename{dyckPath110}
\begin{tikzpicture}
    \draw[gray!30, thin] (0,0) grid (3,3);
    \latticePath{0,0}{nnenee}
\end{tikzpicture}

\tikzsetnextfilename{dyckPath210}
\begin{tikzpicture}
    \draw[gray!30, thin] (0,0) grid (3,3);
    \latticePath{0,0}{nnneee}
\end{tikzpicture}


\end{document}
